\section{Model Comparison and Relation to Turing Machines}
\label{sec:comparison}

The fixed 200-comment test split ranks mean macro F1 as DistilBERT (0.843
$\pm$ 0.006), Dense (0.741 $\pm$ 0.043), LSTM (0.701 $\pm$ 0.021),
SimpleRNN (0.551 $\pm$ 0.047), then the stratified dummy (0.516 $\pm$ 0.025).
The train/test gaps are 0.093 for DistilBERT and 0.249 for Dense. Standard
deviations describe run variation on this split; DistilBERT used three seeds
and the other models used five.

Final-refit CPU times exclude hyperparameter search: Dense, SimpleRNN and LSTM
took about 0.9--1.5 s, while DistilBERT averaged 10.7 min with high variation.
Inference on 1,000 tokenized comments took 47.19 s for DistilBERT and
0.18--0.22 s for the smaller neural models. These hardware-dependent timings
exclude text cleaning and tokenization. Table~\ref{tab:model-cost} reports the
fit, inference and parameter measurements.

For length $T$, embedding width $d$ and hidden width $h$, Dense mean-pooling
costs $O(Td)$, SimpleRNN costs $O(T(dh+h^2))$ sequential operations, and LSTM
has the same order with a larger gate-related constant. DistilBERT's parallel
self-attention costs $O(T^2d)$ per layer. These terms describe sequence
operations, not full wall-clock time.

\begin{table}[t]
\caption{Mean final-fit time, parameter count, and CPU inference time. Fit time is mean $\pm$ sample SD across five runs, or three for DistilBERT; it excludes hyperparameter search.}
\label{tab:model-cost}
\centering\scriptsize
\setlength{\tabcolsep}{2pt}
\begin{tabular}{@{}lrrr@{}}
\toprule
Model & Fit (s) & Params. & Infer./1,000 (s)\\
\midrule
Dummy & $<0.001$ & 0 & 0.0005\\
Dense & $0.86\pm0.06$ & 76,801 & 0.184\\
SimpleRNN & $1.30\pm0.11$ & 81,985 & 0.185\\
LSTM & $1.51\pm0.12$ & 163,873 & 0.221\\
DistilBERT & $644.9\pm424.1$ & 66.96M & 47.193\\
\bottomrule
\end{tabular}
\end{table}

Interpretability varies: Dummy's rule is transparent but text-blind; Dense
averages inspectable token vectors but loses order; RNN/LSTM hide decisions in
learned states. DistilBERT attention is inspectable, not causal.

Turing-completeness results for recurrent networks and Transformers require
idealized assumptions~\cite{siegelmann1995computational,perez2021attention}.
These trained models have fixed parameters, finite precision and bounded inputs,
so every inference terminates. The positive-comment label is a statistical
target, not a formal language; irony and missing context limit this experiment.
